% Source: https://artofproblemsolving.com/community/c3h1344942_rationalizing_denominators
% Stable identity: aops-topic-1344942-q01
Group two terms and use a difference of squares. Multiply numerator and denominator by $\left(1+\sqrt2\right)-\sqrt3$:
\[
\left(1+\sqrt2+\sqrt3\right)\left(1+\sqrt2-\sqrt3\right)=\left(1+\sqrt2\right)^2-3=3+2\sqrt2-3=2\sqrt2 .
\]
So
\[
\dfrac{1}{1+\sqrt2+\sqrt3}=\dfrac{1+\sqrt2-\sqrt3}{2\sqrt2}=\dfrac{\left(1+\sqrt2-\sqrt3\right)\sqrt2}{4}=\boxed{\dfrac{2+\sqrt2-\sqrt6}{4}}.
\]
The same two-step method (pair up terms, then rationalise the remaining single radical) works for any denominator with three square-root terms.
